% Avsnitt 19.1, ur lectures/L19/appendix/a_system_verification.md.

\section{Att verifiera hela systemet}
\label{c19:sec:verify}

Innan någon del av transporten går att lita på måste kontrollern visas fungera som ett \emph{system}.
Det betyder inte en modul mot en attrapp, utan två oberoende noder på en ledning som bråkar om den så
som riktiga noder gör.

\subsection*{Att modellera bussen}
Ingenting i konstruktionen modellerar den delade bussen, och det är medvetet. \code{can_controller}
driver ett par \code{(tx_bus, bus_en)} och läser \code{rx_bus}, vilket är precis vad en riktig
CAN-transceiver presenterar på sin digitala sida (\chapref{c10:ch}). Wired-AND:et bor i testbänken,
och på hårdvara bor det i transceivern.

\file{can_controller_tb.vhd} instansierar två noder på en ledning och kombinerar dem direkt. Varje
nods effektiva utgång är dess drivna värde när \code{bus_en = '1'}, annars släppt (\code{'1'}), och
bussen är deras AND, så den läser dominant i samma ögonblick som \emph{någon} av noderna drar låg:

\begin{vhdlcode}
bus_line <= (a_tx_bus or not a_bus_en) and (b_tx_bus or not b_bus_en);
\end{vhdlcode}

Den enda raden är hela \secref{c10:sec:can}:s ”vilken dominant som helst slår alla recessiva”, och
det är den som gör arbitreringen testbar i simulering.

Värd att räkna igenom snarare än att läsa. Varje nod bidrar med \code{tx_bus or not bus_en}, som är
\code{1} så fort den släppt bussen, och ledningen är de två bidragen AND:ade:

\begin{narrowtable}{llccccl}
  \thead{Nod A} & \thead{Nod B} & \thead{\code{a_tx_bus}} & \thead{\code{a_bus_en}} &
    \thead{\code{b_tx_bus}} & \thead{\code{b_bus_en}} & \thead{\code{bus_line}} \\
  \midrule
  släppt & släppt & \code{-} & \code{0} & \code{-} & \code{0} & \code{1} recessiv \\
  driver dominant & släppt & \code{0} & \code{1} & \code{-} & \code{0} & \code{0} dominant \\
  släppt & driver dominant & \code{-} & \code{0} & \code{0} & \code{1} & \code{0} dominant \\
  driver dominant & driver dominant & \code{0} & \code{1} & \code{0} & \code{1} & \code{0}
    dominant \\
\end{narrowtable}

En släppt nods \code{tx_bus} är \code{-} därför att den faktiskt inte spelar någon roll:
\code{not bus_en} ensamt gör dess bidrag till \code{1}, så en nod som släppt taget kan inte påverka
ledningen oavsett vad dess utgång råkar ligga på. Över alla sexton ingångskombinationer läser
ledningen dominant exakt när minst en nod har \code{bus_en = '1'} med \code{tx_bus = '0'}, vilket är
regeln uttryckt som hårdvara.

Två rader bär arbitreringsargumentet. \textbf{Rad två är själva tvekampen:} A driver dominant, B
sänder recessivt och har därför släppt, och ledningen läser dominant. B, som övervakar medan den
sänder, ser att den sänt recessivt och att bussen läser dominant, och det är precis det förlustvillkor
\chapref{c17:ch} kontrollerar. Samtidigt läser A tillbaka exakt den bit den sände och märker
ingenting.

\textbf{Rad ett är varför förloraren kan gå.} När B slutat driva bidrar den med \code{1} för alltid,
så ledningen är A:s ensam. Förlorarens reträtt är osynlig för vinnaren, och det är det som gör en
kollision till en icke-händelse i stället för något båda noderna måste återhämta sig från.

Rad fyra är värd en blick också: när båda noderna driver dominant är ledningen dominant och var och
en läser tillbaka det den sände, så ingen av dem kan märka att den andra finns. Det är det normala
tillståndet för varje dominant bit i identifieraren innan de två ID:na skiljer sig åt.

\subsection*{De två scenarierna}
\begin{enumerate}
  \item \textbf{Mottagning.} Nod A sänder en ram på 5 bytes (ID \code{0x123}, data
    \code{11 22 F8 44 55}) medan nod B tar emot den. \code{0xF8} framtvingar en stoppbit inne i
    datafältet, så båda sidors CRC-motorer måste hoppa över den. Testbänken kontrollerar A:s
    \code{tx_done} och B:s \code{rx_valid}, att B rekonstruerade ID:t, DLC:n och \textbf{varje}
    databyte på sin fasta position, och att ingen av noderna höjde \code{error}.
  \item \textbf{Arbitrering.} Noderna A (ID \code{0x000}) och B (ID \code{0x001}) begär båda under
    samma cykel, och skiljer sig bara i identifierarens sista bit. A måste vinna och fullborda
    normalt, ovetande om att någon tvekamp ägt rum; B måste upptäcka förlusten vid den biten, höja
    \code{error} och sluta driva. Kontrollen körs efter att A:s ram är klar, så den bekräftar också
    att B:s \code{error} fortfarande går att läsa då, i stället för att ha suddats av ett återinträde
    i \code{STATE_START}.
\end{enumerate}

En sändande nod tar aldrig emot sin egen ram (\code{rx_shift_reg} går bara medan \code{role = '0'}),
så båda scenarierna behöver verkligen två noder; inget av dem skulle gå att testa med en ensam nod.

De två encykelspulser det första scenariot väntar på, A:s \code{tx_done} och B:s \code{rx_valid},
\textbf{låses} i testbänken, eftersom två oberoende tillståndsmaskiner inte behöver pulsa på samma
flank, och ett rakt \code{wait until (a = '1' or b = '1')} skulle återuppta på den första och missa
den andra.

\subsection*{Vad testbänken inte täcker}
\label{c19:sub:gaps-tb}

Två godkända scenarier är inte samma sak som en bevisad konstruktion, och att veta var en testbänk
slutar hör till att läsa den. Fem luckor är värda att namnge, eftersom var och en är ett ställe där
den här konstruktionen antingen är förenklad med flit eller helt enkelt otestad:
\begin{itemize}
  \item \textbf{Båda noderna delar en \code{clock} och en \code{reset_n}.} Riktiga noder går på
    oberoende oscillatorer och driver isär. Var precis med vad det lämnar otestat, för det är inte
    \code{resync} självt: båda noderna pulsar \code{resync} vid varje rams början (\chapref{c16:ch}),
    och den pulsen är det som fasriktar B:s sampelpunkter mot det godtyckliga ögonblick då A började
    sända, så fall 1 misslyckas utan den. Vad en delad klocka aldrig kan pröva är \textbf{drift}, två
    timrar som startar i fas och vandrar isär mitt i ramen, fallet som Övning~\ref{c12:ex:drift}
    kvantifierar och skälet till att riktig CAN synkroniserar om vid varje flank från recessiv till
    dominant i stället för en gång per ram. Det här är den största luckan, och att sluta den i
    simulering kräver en andra klockgenerator med en avsiktlig frekvensavvikelse.
    \Secref{c19:sec:bringup} får en ram tvärs över \textbf{två kort} med två separata kristaller,
    äntligen genuint oberoende klockor, med ett ärligt förbehåll noterat där.
  \item \textbf{Ingen ram förvanskas någonsin.} Ingenting matar in ett stoppfel, vänder en bit på
    ledningen eller skadar en CRC, så mottagarsidans \code{error}-vägar prövas bara på modulnivå av
    \code{rx_shift_reg_tb} och \code{crc15_tb}, aldrig hela vägen igenom.
  \item \textbf{ACK-luckan drivs men läses aldrig tillbaka.} En mottagare med giltig CRC drar den
    dominant (\chapref{c17:ch}), men sändaren kontrollerar den aldrig, så en ram som ingen kvitterat
    fullbordas precis som en som blev kvitterad. Det är en medveten förenkling
    (\secref{c19:sec:gaps}), och testbänken ärver den i stället för att fånga den.
  \item \textbf{Arbitreringen är alltid tvåvägs, och förloras alltid vid samma bit.} \code{0x000} mot
    \code{0x001} skiljer sig i identifierarens sista bit, så förlusten upptäcks i samma tillstånd
    varje körning. Övning~\ref{c19:ex:thirdnode} lägger till en tredje nod och flyttar den förlorande
    biten tidigare.
  \item \textbf{Ingen nod sänder någonsin igen efter ett avbrott.} Nod B förlorar arbitreringen i
    fall 2 och blir aldrig ombedd att skicka en ram till, så ingenting här kontrollerar att en nod
    över huvud taget återhämtar sig från ett avbrott. Det spelar större roll än det låter:
    \code{can_controller} nollställer \code{crc15} i \code{STATE_START} (\chapref{c16:ch}) just för
    att en övergiven ram inte ska kunna lämna sina rester efter sig, och den här testbänken skulle
    passera precis lika glatt med den nollställningen borttagen. Övning~\ref{c19:ex:afterabort} är
    där du får reda på vad den gjorde.
\end{itemize}
